\documentclass[11pt,a4paper]{letter}
\usepackage[utf8]{inputenc}
\usepackage[T1]{fontenc}
\usepackage[french]{babel}
\usepackage[left=2.5cm,right=2.5cm,top=2cm,bottom=2cm]{geometry}
% \usepackage{hyperref} % Désactivé pour compatibilité portail recrutement
\usepackage{xcolor}

% --- Couleurs Eviden ---
\definecolor{primary}{RGB}{0, 55, 110}

\signature{Loïc Aron MBASSI EWOLO}
\address{Loïc Aron MBASSI EWOLO\\
Cergy, France\\
(+33) 7 59 52 33 98\\
loicaron.mbassiewolo@ensea.fr\\
LinkedIn : loic-mbassi | GitHub : Nameless0l}

\begin{document}

\begin{letter}{Eviden / Atos\\
Pôle Radio Logicielle\\
Aix-en-Provence}

\opening{Madame, Monsieur,}

Élève-ingénieur en 2A à l'ENSEA (double diplôme avec Polytechnique Yaoundé), je candidate au stage Firmware FPGA sur plateforme Xilinx RFSoC.

Je travaille actuellement sur un projet FPGA à l'ENSEA : intégration d'un soft-processeur Nios V sur Cyclone V avec Platform Designer. J'ai l'habitude du flow Quartus (synthèse, timing, simulation ModelSim) et de l'interconnexion de blocs via bus Avalon. Je n'ai pas encore travaillé sur Xilinx/Vivado, mais la logique reste la même et je suis motivé pour m'adapter.

Côté traitement du signal, je suis les cours à l'ENSEA (filtrage, FFT, acquisition) et j'ai travaillé sur du traitement de signaux ECG et de données capteurs. La partie acquisition multi-voies synchrones du sujet m'intéresse particulièrement.

En parallèle, j'ai une base solide en \textbf{C/C++} (projet système multiprocessus, STM32) et \textbf{Python}. Je travaille sous Linux au quotidien.

Le contexte Défense et radio logicielle m'attire -- c'est un domaine où la rigueur technique compte vraiment. Et le fait que le bureau d'études ait une composante R\&D forte correspond à ce que je recherche.

Disponible pour un entretien.

\closing{Cordialement,}

\end{letter}
\end{document}
